\begin{ex}
Mostrare che l'insieme $X=\{p\in\mathbb N\mid p\text{ è primo}\}$ è ricorsivo.
\end{ex}

\begin{sol}
Dobbiamo mostrare che la funzione caratteristica $\mathbb 1_X$ è ricorsiva.
\ 
\begin{enumerate}
\item Costruiamo la funzione \texttt{r} che calcola il resto della divisione. Sia
\[
\texttt{r}\colon\mathbb N^2\to\mathbb N,\qquad
(x,y)\longmapsto
\begin{cases}
[x]_y & y>0,\\
\bot & y=0.
\end{cases}
\]
Mostriamo che è ricorsiva. Se $y>0$ si ha $x=ky+r$ con $0\le r<y$, quindi dobbiamo trovare $k$. Sia allora
\[
\texttt{q}\colon\mathbb N^2\to\mathbb N,\qquad
(x,y)\longmapsto
\begin{cases}
k & k\text{ minimo tale che }ky>x,\\
\bot & \text{altrimenti.}
\end{cases}
\]
Osserviamo che se $y=0$ non esiste alcun $k$ tale che $ky>x$, e quindi $\texttt{q}(x,0)=\bot$; vale anche il viceversa, dunque $\texttt{q}$ è definita se e solo se $y\neq0$. Per costruirla poniamo $\texttt{q}=\min(\texttt{sigma})$, dove
\[
\texttt{sigma}\colon\mathbb N^3\to\mathbb N,\qquad (x,y,i)\longmapsto 1-\texttt{lesser}(x,iy).
\]
In questo modo, per ogni $(x,y)\in\mathbb N\times(\mathbb N\setminus\{0\})$ esiste $r$ con $0\le r<y$ tale che $x=(\texttt{q}(x,y)-1)\,y+r$ (si noti che $\texttt{q}(x,y)\ge1$, perché $0\cdot y=0\not>x$). Possiamo quindi definire $\texttt{r}(x,y):=x-y\,(\texttt{q}(x,y)-1)$. Dalla definizione di $\texttt{q}$ abbiamo
\[
\texttt{r}\colon\mathbb N^2\to\mathbb N,\qquad
(x,y)\longmapsto
\begin{cases}
x-y\,(\texttt{q}(x,y)-1) & y\neq0,\\
\bot & y=0,
\end{cases}
\]
che è una funzione ricorsiva.

\item Vogliamo costruire $\mathbb 1_X$. I casi base $0$ e $1$ li consideriamo a parte, cioè
\[
\mathbb 1_X=(1-\mathbb 1_{\{0\}})(1-\mathbb 1_{\{1\}})\cdot\texttt{tau},
\]
dove $\texttt{tau}$ deve decidere i numeri primi. Poniamo
\[
\texttt{test}_{=}(x,y):=\Delta\big(1,\Delta(x,y)+\Delta(y,x)\big),
\]
che vale $1$ se $x=y$ e $0$ altrimenti. Per costruire $\texttt{tau}$ definiamo
\[
\texttt{f}\colon\mathbb N\to\mathbb N,\qquad
x\longmapsto k+2\quad\text{con $k$ minimo tale che }[x]_{k+2}=0\ \text{ oppure }\ x=1.
\]
Un tale $k$ esiste sempre (per $x=0$ si ha $k=0$, per $x=1$ vale la seconda condizione, per $x\ge2$ si può prendere $k=x-2$), quindi $\texttt{f}$ è totale. Si ha $\texttt{f}(0)=\texttt{f}(1)=2$ e, per $x\ge2$, $\texttt{f}(x)$ è il più piccolo divisore di $x$ maggiore o uguale a $2$.

Questa funzione è ricorsiva, infatti sia
\[
\texttt{eta}\colon\mathbb N^2\to\mathbb N,\qquad
(x,i)\longmapsto
\begin{cases}
0 & \text{se } [x]_{i+2}=0 \text{ oppure } x=1,\\
1 & \text{altrimenti,}
\end{cases}
\]
e abbiamo $\texttt{f}(x)=\min(\texttt{eta})(x)+2$.

Ora, per $x\ge2$, $x$ è primo se e solo se $\texttt{f}(x)=x$, e quindi
\[
\texttt{tau}(x):=\texttt{test}_{=}\big(\texttt{f}(x),x\big).
\]
In conclusione
\[
\mathbb 1_X=(1-\mathbb 1_{\{0\}})(1-\mathbb 1_{\{1\}})\cdot\texttt{tau}.
\]

\item Troviamo ora l'espressione per ogni pezzo. Sia
\[
\begin{aligned}
\mathsf{SIGMA}:=\mathsf{COMP}\Big(\mathsf{CONCAT}\big(\mathsf C_{3,1},\ &\mathsf{COMP}\big(\mathsf{CONCAT}\big(\mathsf P_{3,1},\\
&\mathsf{COMP}(\mathsf{CONCAT}(\mathsf P_{3,3},\mathsf P_{3,2}),\mathsf{MULT})\big),\\
&\mathsf{LESSER}\big)\big),\mathsf{DIFF}\Big),
\end{aligned}
\]
e vale $\llbracket\mathsf{SIGMA}\rrbracket=\texttt{sigma}$. Quindi
\[
\mathsf Q:=\mathsf{MIN}(\mathsf{SIGMA}),\qquad \llbracket\mathsf Q\rrbracket=\texttt{q},
\]
e di nuovo
\[
\begin{aligned}
\mathsf R:=\mathsf{COMP}\Big(\mathsf{CONCAT}\big(\mathsf P_{2,1},\ &\mathsf{COMP}\big(\mathsf{CONCAT}\big(\mathsf P_{2,2},\\
&\mathsf{COMP}(\mathsf{CONCAT}(\mathsf Q,\mathsf C_{2,1}),\mathsf{DIFF})\big),\mathsf{MULT}\big)\big),\mathsf{DIFF}\Big),
\end{aligned}
\]
e abbiamo $\llbracket\mathsf R\rrbracket=\texttt{r}$. Inoltre
\[
\begin{aligned}
\mathsf{TEST}_{=}:=\mathsf{COMP}\Big(\mathsf{CONCAT}\big(\mathsf C_{2,1},\ &\mathsf{COMP}\big(\mathsf{CONCAT}\big(\mathsf{DIFF},\\
&\mathsf{COMP}(\mathsf{CONCAT}(\mathsf P_{2,2},\mathsf P_{2,1}),\mathsf{DIFF})\big),\mathsf{SUM}\big)\big),\mathsf{DIFF}\Big),
\end{aligned}
\]
e vale $\llbracket\mathsf{TEST}_{=}\rrbracket=\texttt{test}_{=}$. Poniamo poi
\[
\begin{aligned}
\mathsf{ISZERO}&:=\mathsf{COMP}\big(\mathsf{CONCAT}(\mathsf P_{1,1},\mathsf C_{1,0}),\mathsf{TEST}_{=}\big),\\
\mathsf{ISUNO}&:=\mathsf{COMP}\big(\mathsf{CONCAT}(\mathsf P_{1,1},\mathsf C_{1,1}),\mathsf{TEST}_{=}\big),
\end{aligned}
\]
e quindi
\[
\begin{aligned}
\mathsf{NOTZERO}&:=\mathsf{COMP}\big(\mathsf{CONCAT}(\mathsf C_{1,1},\mathsf{ISZERO}),\mathsf{DIFF}\big),\\
\mathsf{NOTUNO}&:=\mathsf{COMP}\big(\mathsf{CONCAT}(\mathsf C_{1,1},\mathsf{ISUNO}),\mathsf{DIFF}\big),
\end{aligned}
\]
e abbiamo $\llbracket\mathsf{NOTZERO}\rrbracket=1-\mathbb 1_{\{0\}}$ e $\llbracket\mathsf{NOTUNO}\rrbracket=1-\mathbb 1_{\{1\}}$.

Infine $\texttt{eta}(x,i)=\texttt{a}(x,i)\cdot\texttt{b}(x,i)$, dove
\[
\texttt{a}(x,i)=1-\mathbb 1_{\{1\}}(x),\qquad
\texttt{b}(x,i)=1-\texttt{test}_{=}\big(0,[x]_{i+2}\big).
\]
Se scriviamo
\[
\mathsf A:=\mathsf{COMP}\big(\mathsf{CONCAT}\big(\mathsf C_{2,1},\ \mathsf{COMP}(\mathsf P_{2,1},\mathsf{ISUNO})\big),\mathsf{DIFF}\big),
\]
\[
\begin{aligned}
\mathsf B:=\mathsf{COMP}\Big(\mathsf{CONCAT}\Big(\mathsf C_{2,1},\ &\mathsf{COMP}\big(\mathsf{CONCAT}\big(\mathsf C_{2,0},\\
&\mathsf{COMP}\big(\mathsf{CONCAT}\big(\mathsf P_{2,1},\\
&\mathsf{COMP}(\mathsf{CONCAT}(\mathsf P_{2,2},\mathsf C_{2,2}),\mathsf{SUM})\big),\mathsf R\big)\big),\\
&\mathsf{TEST}_{=}\big)\Big),\mathsf{DIFF}\Big),
\end{aligned}
\]
e quindi
\[
\mathsf{ETA}:=\mathsf{COMP}\big(\mathsf{CONCAT}(\mathsf A,\mathsf B),\mathsf{MULT}\big),\qquad \llbracket\mathsf{ETA}\rrbracket=\texttt{eta},
\]
e quindi
\[
\mathsf F:=\mathsf{COMP}\big(\mathsf{CONCAT}(\mathsf C_{1,2},\mathsf{MIN}(\mathsf{ETA})),\mathsf{SUM}\big),\qquad \llbracket\mathsf F\rrbracket=\texttt{f},
\]
\[
\mathsf{TAU}:=\mathsf{COMP}\big(\mathsf{CONCAT}(\mathsf F,\mathsf P_{1,1}),\mathsf{TEST}_{=}\big),\qquad \llbracket\mathsf{TAU}\rrbracket=\texttt{tau},
\]
e quindi, indicando con $\mathsf{MULT}^3$ il prodotto di tre fattori,
\begin{gather*}
\mathsf{IND}_X:=\mathsf{COMP}\big(\mathsf{CONCAT}^3(\mathsf{NOTZERO},\mathsf{NOTUNO},\mathsf{TAU}),\mathsf{MULT}^3\big),\\
\llbracket\mathsf{IND}_X\rrbracket=\mathbb 1_X.
\end{gather*}
Dunque $\mathbb 1_X$ è ricorsiva e $X$ è un insieme ricorsivo.
\end{enumerate}
\end{sol}

\newpage